\documentclass[../main.tex]{subfiles}
\begin{document}

\chapter*{Overview at a Glance}
\addcontentsline{toc}{chapter}{Overview at a Glance}
\markboth{Overview at a Glance}{}
\label{ch:overview}

This half of the guide is the fast map. It collects the snapshot card and the part summaries for all six weeks, with no detailed content in between, so the whole shape of the course can be scanned in a few pages. Each part summary ends with a link to the section in the detailed reference half where its bullets and prose live. The detailed chapters in turn link back to the matching week here.

% ===================== WEEK 1 =====================
\section*{Week 1 — Introduction to Research}
\phantomsection\label{ov:week1}

\begin{figure}[H]
  \centering
  \includegraphics[width=\linewidth]{week1.png}
\end{figure}

\begin{snapcard}{conceptgreen}{Week 1 — Introduction to Research}{Course start, no prerequisites}{Conceptual and technical design (\autoref{ch:week2}), sampling (\autoref{ch:week3}), reports (\autoref{ch:week6})}
\begin{itemize}
  \item The research process runs in six fixed steps: the problem, the sample, the design, measure, analyse, conclude.
  \item A hypothesis is a concrete prediction, and its form fixes whether the study is descriptive, relational, or causal.
  \item The variable is the smallest unit of research language, and its values are attributes that must be exhaustive and mutually exclusive.
  \item An observed relationship may be correlational, causal, or produced by a third variable, so correlation never proves causation.
  \item Deduction reasons from general to specific and induction from specific to general, and the hourglass joins the two in one study.
  \item Feasibility rather than interest decides whether an idea becomes a project, and the literature review sets that idea in context.
\end{itemize}
\end{snapcard}

\begin{partbox}\textbf{Part 1 — The Research Process.} Six-step pipeline: problem → sample → design → measure → analyse → conclude. Each step constrains the next. Design is expanded in Week~2 (\autoref{ch:week2}); measure in Week~3 (\autoref{ch:week3}). Full detail in \autoref{sec:research-process}.\end{partbox}

\begin{partbox}\textbf{Part 2 — The Language of Research.} Defines hypothesis, study type (descriptive / relational / causal), variable, attribute, and relationship nature. Correlation is not causation. Full detail in \autoref{sec:language-research}.\end{partbox}

\begin{partbox}\textbf{Part 3 — The Hourglass: Deduction and Induction.} Deduction narrows from broad question to measurement; induction widens back to generalised claim. Induction is fallible. Operationalisation at the neck is formalised in Week~2 (\autoref{sec:operationalization}). Full detail in \autoref{sec:hourglass}.\end{partbox}

\begin{partbox}\textbf{Part 4 — Where Research Ideas Come From.} Ideas arise from field problems, knowledge gaps, and prior conclusions. Feasibility (time, ethics, cost, access) is a gate. The literature review locates the gap that becomes the Week~2 objective (\autoref{sec:objectives-rqs}). Full detail in \autoref{sec:research-ideas}.\end{partbox}

% ===================== WEEK 2 =====================
\section*{Week 2 — Research Design}
\phantomsection\label{ov:week2}

\begin{figure}[H]
  \centering
  \includegraphics[width=\linewidth]{week2.png}
\end{figure}

\begin{snapcard}{conceptgreen}{Week 2 — Research Design}{Language of research, induction and deduction, steps of a research project (\autoref{ch:week1})}{Reference list (\autoref{ch:week6}), project management (\autoref{ch:week5})}
\begin{itemize}
  \item Research design splits into conceptual design, which fixes what, why, and how much, and technical design, which fixes how, where, and when.
  \item Conceptual design has four phases: the objective, the research framework, the research question or conceptual model, and definition with operationalisation.
  \item A good research objective is informative, useful, feasible, and clear, and it names a contribution and the way that contribution is delivered.
  \item Research questions must be open-ended, number at least two, and be relevant, researchable, anchored, and precise.
  \item Technical design has three phases: research strategy, research material, and research planning.
  \item Reproducibility and the code of conduct govern conduct, and any work with human subjects requires HREC approval.
\end{itemize}
\end{snapcard}

\begin{partbox}\textbf{Part 1 — Strategic Areas.} Conceptual design (what, why, how much — 4 phases) must precede technical design (how, where, when — 3 phases). Technical choices are undefined until conceptual ones are settled. Full detail in \autoref{sec:conceptual-technical}.\end{partbox}

\begin{partbox}\textbf{Part 2 — Conceptual Design.} Builds the objective, research questions, framework, and operationalisation. A weak question set is the usual reason a technical design cannot be specified. Checklists in Appendix~\ref{ch:checklists}. Full detail in \autoref{sec:objectives-rqs}.\end{partbox}

\begin{partbox}\textbf{Part 3 — Technical Design.} Turns operationalised concepts into strategy, materials, and schedule. Planning tools (WBS, Gantt) are developed in Week~5 (\autoref{sec:project-mgmt}). Full detail in \autoref{sec:strategy-materials}.\end{partbox}

\begin{partbox}\textbf{Part 4 — Ethics and Conduct.} Reproducibility, traceability, and proper credit govern all data and method choices. Human-subject work requires HREC approval before proceeding. Same traceability requirement as EOP and the DMP (\autoref{sec:eop}, \autoref{sec:dmp}). Full detail in \autoref{sec:ethics}.\end{partbox}

% ===================== WEEK 3 =====================
\section*{Week 3 — Research Material and Data}
\phantomsection\label{ov:week3}

\begin{figure}[H]
  \centering
  \includegraphics[width=\linewidth]{week3.png}
\end{figure}

\begin{snapcard}{empiricred}{Week 3 — Research Material and Data}{Steps of a research project (\autoref{ch:week1})}{Verification and validation, sustainable research (\autoref{ch:week4}), data management (\autoref{ch:week5})}
\begin{itemize}
  \item The population is the entire collection in the field, and the sample is the representative subset that is actually studied.
  \item The sampling method narrows from the theoretical population, through the study population and the sampling frame, to the sample.
  \item Probability sampling supports generalisation to the population, while non-probability sampling never can and must say so explicitly.
  \item The standard error measures the precision of a sample, falling as the sample size grows and as the standard deviation shrinks.
  \item Data is qualitative or quantitative, and its type follows from the strategy, the materials, and the defined variables.
  \item Data quality must be checked before analysis, and a detailed research logbook is what makes the whole project traceable.
\end{itemize}
\end{snapcard}

\begin{partbox}\textbf{Part 1 — Population, Sample, and the Sampling Method.} Population = full field; sample = studied subset. Representativeness licenses generalisation; without it conclusions hold only for the sample. Four-stage funnel: theoretical population → study population → sampling frame → sample. Full detail in \autoref{sec:sampling}.\end{partbox}

\begin{partbox}\textbf{Part 2 — Types of Samples and Standard Error.} Probability sampling supports generalisation; non-probability cannot and must state this. Standard error quantifies precision: larger spread increases it, larger sample reduces it (\autoref{sec:eop}). Full detail in \autoref{sec:sample-types}.\end{partbox}

\begin{partbox}\textbf{Part 3 — Data Sources, Quality, and the Logbook.} Covers data sources, GDPR and knowledge-embargo constraints, qualitative/quantitative split, quality checks, and the logbook. Personal data and protected fields are hard gates that can stop a project. Logbook is formalised into the DMP in Week~5 (\autoref{sec:dmp}). Full detail in \autoref{sec:data-sources}.\end{partbox}

% ===================== WEEK 4 =====================
\section*{Week 4 — Validation and Sustainability}
\phantomsection\label{ov:week4}

\begin{figure}[H]
  \centering
  \includegraphics[width=\linewidth]{week4.png}
\end{figure}

\begin{snapcard}{empiricred}{Week 4 — Validation and Sustainability}{Sampling and data (\autoref{ch:week3})}{Data management and project management (\autoref{ch:week5})}
\begin{itemize}
  \item Verification proves the method was used correctly and suits the topic, while validation proves the outcomes are true on strong evidence.
  \item Both are needed for credibility, professional use, and reproducibility, and both must be designed into the plan with time reserved.
  \item In industry, verification is internal quality control against specifications and validation is external quality assurance of fitness for purpose.
  \item The Engineering for One Planet framework weighs research on its value, its reproducibility, and its resource management.
  \item Reproducibility depends on good data management and is blocked by complexity, technology change, human error, and intellectual property.
  \item Resource management spans digital habits and methodology choices, each judged on whether the cost matches the accuracy actually needed.
\end{itemize}
\end{snapcard}

\begin{partbox}\textbf{Part 1 — Verification and Validation.} Verification: method applied correctly and suited to the topic. Validation: outcomes are true on strong evidence. Both must be scheduled — never assume time will be available (\autoref{sec:strategy-materials}). Full detail in \autoref{sec:vv}.\end{partbox}

\begin{partbox}\textbf{Part 2 — Engineering for One Planet.} Three axes: value (near- and long-term benefit/harm), reproducibility (barriers: complexity, tech change, error, IP), resource management (digital habits and methodology choices). Data management addresses the reproducibility barriers (\autoref{sec:dmp}); resource axis ties back to sample-size trade-off (\autoref{sec:sample-types}). Full detail in \autoref{sec:eop}.\end{partbox}

% ===================== WEEK 5 =====================
\section*{Week 5 — Data Management and Project Planning}
\phantomsection\label{ov:week5}

\begin{figure}[H]
  \centering
  \includegraphics[width=\linewidth]{week5.png}
\end{figure}

\begin{snapcard}{integblue}{Week 5 — Data Management and Project Planning}{Technical design (\autoref{ch:week2}), sustainable research and data (\autoref{ch:week3}, \autoref{ch:week4})}{Reporting (\autoref{ch:week6})}
\begin{itemize}
  \item A data management plan has six sections, from data description to responsibilities, and must be designed before any data is collected.
  \item FAIR data is findable, accessible, interoperable, and reusable, which is the goal that good documentation and metadata serve.
  \item The 3-2-1 rule keeps three copies on two media with one off-site, fixing the backup strategy from the start.
  \item Project management gives a research project structured timelines, optimised resources, early risk detection, and clear communication.
  \item A Work Breakdown Structure divides the project into work packages and tasks down to a level where each task can be estimated.
  \item A Gantt chart schedules the work and exposes the critical path, the chain that sets the shortest possible project duration.
\end{itemize}
\end{snapcard}

\begin{partbox}\textbf{Part 1 — The Data Management Plan.} Six-section DMP must exist before data collection. Encodes the 3-2-1 backup rule and licensing choices. Concrete form of the Week~4 reproducibility duty (\autoref{sec:eop}). Full detail in \autoref{sec:dmp}.\end{partbox}

\begin{partbox}\textbf{Part 2 — Project Management in Research.} Develops the Week~2 planning step (\autoref{sec:strategy-materials}) into task estimation, WBS, and Gantt chart. The critical path fixes the minimum project duration. Planning only assembles once both designs are complete. Full detail in \autoref{sec:project-mgmt}.\end{partbox}

% ===================== WEEK 6 =====================
\section*{Week 6 — Reporting and Writing}
\phantomsection\label{ov:week6}

\begin{figure}[H]
  \centering
  \includegraphics[width=\linewidth]{week6.png}
\end{figure}

\begin{snapcard}{integblue}{Week 6 — Reporting and Writing}{Steps of a research project (\autoref{ch:week1}), conceptual design (\autoref{ch:week2})}{The Research Methodologies report (\autoref{ch:reportguide})}
\begin{itemize}
  \item A technical report conveys specific information about a technical subject to a specific audience for a specific and practical purpose.
  \item The Research Methodologies report conveys a research plan to a technical peer in order to organise the research activities.
  \item Every report has a head that introduces, a main body of connected sections, and feet that conclude without adding new information.
  \item Choose one reference style and keep it, either author-date such as APA or numbered such as AIAA, each with its own trade-off.
  \item A reference list should use at least three document types, lean on peer-reviewed articles, and stay mostly under ten years old.
  \item Reader-centred writing guides the reader in through titles, figures, and conclusions first, and a final checklist closes the report.
\end{itemize}
\end{snapcard}

\begin{partbox}\textbf{Part 1 — Technical Reports.} Defines the technical report (specific subject, audience, purpose) and maps the head/body/feet skeleton onto the RM report. Section-by-section mapping in \autoref{ch:reportguide}. Full detail in \autoref{sec:reports}.\end{partbox}

\begin{partbox}\textbf{Part 2 — Referencing.} Source quality rules (types, currency, peer-review ratio) and the two citation styles. Author-date identifies the source in-text; numbered keeps the text clean. Mixing styles is an error. Full detail in \autoref{sec:referencing}.\end{partbox}

\begin{partbox}\textbf{Part 3 — Writing Style.} Define the target reader first. Titles, figures, and conclusions are read first and must stand alone. 4Cs + graph conventions + AI cautions + final checklist. Full detail in \autoref{sec:writing-style}.\end{partbox}

\end{document}
